\chapter{Literature Review}
\label{Literature_Review}

Five bodies of work bear on the proposed system: learning management systems, mind mapping, gamification, learning analytics dashboards, and retrieval-augmented assistants. The review below treats each in turn, and closes with the gap that the five have in common.

\section{Learning Management Systems}

Moodle is the reference case for an open-source platform built on a pedagogical commitment rather than on a content model. Dougiamas and Taylor \cite{dougiamas2004moodle} describe a system in which an instructor establishes a space and students construct knowledge inside it through discussion, and that principle shaped both the interface and the data model. Twenty years of adoption later, the default course page is still a list of resources with completion ticks beside them, and the gradebook is still the instrument through which progress is read.

Cloud-based tools in blended learning were surveyed by Al-Samarraie and Saeed \cite{alsamarraie2018systematic}, who classified 29 studies into synchronous tools, learning management systems and social networking tools. Collaboration was well supported in all three categories, through shared files, co-editing and discussion threads. What the survey did not address is the student's own estimate of their position in a course, which is the question this work begins from.

Generative artificial intelligence entered the same conversation more recently. Bozkurt et al. \cite{bozkurt2023speculative} assembled reflections from more than thirty education researchers using a speculative method, and two of their observations bear directly on the present design. The division of labour between a teacher and an automated assistant has to be decided rather than allowed to emerge, and students are poor judges of whether a fluent answer is a reliable one. The second observation is the reason retrieval here is confined to course material.

\section{Mind Mapping in Education}

The technique that Buzan and Buzan \cite{buzan2006mind} popularised rests on a correspondence between the shape of a drawing and the shape of recall: a central concept, branches radiating outward, colour and image carrying part of the meaning. Whatever the strength of the underlying claim, the drawing is easier to inspect than a list, which is why the technique survived as a study aid long after its theoretical framing was questioned.

Format matters to the task, and Eppler \cite{eppler2006comparison} established where each format belongs by comparing concept maps, mind maps, conceptual diagrams and visual metaphors. A mind map suits a hierarchical, tree-shaped body of knowledge, which a syllabus typically is. A concept map can express a relationship that a tree cannot, such as a link between two topics on different branches. The design in Chapter 5 keeps a tree for the syllabus and permits additional edges between topics, which is a compromise between the two.

That mind mapping remained a private study technique is the striking part of this literature. In two decades of learning management systems, none of the widely adopted platforms has made a map the primary surface for navigation and progress. The gap is not conceptual; it is an absence of implementation.

\section{Gamification in Learning}

Deterding et al. \cite{deterding2011game} supplied the definition that most subsequent work uses, and the distinction that matters here: a complete game is not the same thing as the selective use of one game mechanic. Synapse borrows a mechanic, the skill tree, without presenting a course as a game, and the distinction is what keeps the borrowing defensible.

Evidence on outcomes is favourable but conditional. Hamari et al. \cite{hamari2014does} reviewed empirical studies across domains and found predominantly positive effects on engagement, with two qualifications worth recording: effects were stronger in education and health than in commerce, and they weakened when the game layer was disconnected from the underlying task. A progression indicator driven by coursework is not exposed to the second effect, since the feedback is a by-product of the work.

Coverage of specific mechanics comes from Dicheva et al. \cite{dicheva2015gamification}, whose systematic mapping classified educational gamification by mechanic, level, subject and evaluation method. Points, badges and leaderboards dominate, and the call was for controlled experiments measuring learning rather than engagement counts. Visual progression of the skill-tree kind is thin in that map, which is an argument for documenting a working instance.

Two commercial products are usually cited as successes, and both are subject-specific: Duolingo for language practice and Khan Academy for mathematics. A general-purpose platform that applies tree-shaped progression to an arbitrary syllabus has not been attempted at scale.

\section{Learning Analytics and Visual Dashboards}

The discipline's own history is instructive. Siemens \cite{siemens2013learning} traced its formation from scattered publications into conferences, centres and a research society, and argued that counting clicks is not the same as modelling learning. Privacy and consent were raised in the same early work, before dashboards became routine.

What dashboards actually display was surveyed by Verbert et al. \cite{verbert2014learning}, who grouped implementations by setting: lectures, group work, and awareness or reflection in online learning. Most aimed at awareness and goal setting, and most reached for basic charts. The recommendation was for visualisations beyond bars and lines, which describes precisely what is missing.

The largest review of the area, by Schwendimann et al. \cite{schwendimann2017perceiving}, covered 55 studies and found a field still producing prototypes rather than production systems. Network graphs, trees and mindmap layouts were almost absent from the corpus. A second finding is methodological: dashboards were built on assumptions about useful information rather than co-designed with the students and teachers who would read them.

Dual tracking is the specific absence. Across those reviews, no dashboard places the teacher's position in the syllabus beside each student's position, and that is the comparison that identifies a gap while a course is still running.

\section{Artificial Intelligence and Retrieval-Augmented Generation}

Grounding a language model in documents rather than parameters is the contribution of Lewis et al. \cite{lewis2020retrieval}. A retriever selects passages, a generator conditions on them, and the answer acquires a provenance that a parametric model cannot provide. On open-domain question answering the method outperformed a purely parametric baseline, with the gain concentrated in questions whose answers lie outside the training data.

For a course, the mapping is direct: a question is answered from the notes and readings supplied for that subject, not from what the model absorbed during training. Integration into a complete learning management system is where existing work thins out, since most retrieval-augmented tutors are research prototypes operating on a folder of documents. The implementation described here builds the index from material attached to a subject and restricts retrieval to that subject at query time.

Bozkurt et al. \cite{bozkurt2023speculative} add the caution that shapes the guard: a confident wrong answer is worse than no answer when the student cannot check it. Restricting retrieval is the cheapest safeguard available, and it is the one adopted.

\section{Summary of the Literature}

Table \ref{tab:literature_summary} condenses the five areas into current state and identified gap.

\begin{table}[H]
\centering
\caption{Summary of Literature Review Areas and Identified Gaps}
\label{tab:literature_summary}
\resizebox{\textwidth}{!}{
\begin{tabular}{|p{3.5cm}|p{5cm}|p{5cm}|}
\hline
\thead{Area} & \thead{Current State} & \thead{Gap / Opportunity} \\ \hline
LMS Platforms & Content delivery and grade tracking & No visual progress mapping \\ \hline
Mind Mapping & Used as a study tool, not in LMS & Not integrated into progress tracking \\ \hline
Gamification & Effective in specific apps (Duolingo) & Not applied in general-purpose LMS \\ \hline
Learning Analytics & Basic charts and dashboards & No dual tracking (teacher vs student) \\ \hline
AI in Education & RAG reduces hallucinations & Not integrated into LMS with scope control \\ \hline
\end{tabular}}
\end{table}
